% Compile twice: lualatex 20001241.tex
\documentclass[11pt]{article}
\usepackage[margin=25mm]{geometry}
\usepackage{fontspec,unicode-math}
\setmainfont{Latin Modern Roman}
\setmathfont{Latin Modern Math}
\usepackage{amsmath,amssymb,mathtools}
\usepackage{xurl,hyperref}
\hypersetup{colorlinks=true,urlcolor=blue,linkcolor=blue}
\setcounter{secnumdepth}{0}
\setlength{\parindent}{0pt}
\setlength{\parskip}{5pt}
\setlength{\emergencystretch}{3em}
\newcommand{\sref}[2]{\hyperlink{source:#1:#2}{[S#2]}}
\newcommand{\cataloguescope}[1]{\paragraph{Source question.}#1}
\newcommand{\floor}[1]{\left\lfloor#1\right\rfloor}
\begin{document}
% UnsolvedMath ID 20001241; source catalogue code AIM-COMPUTATION-0079
\setcounter{section}{20001240}
\section{Removing term order without hiding the extrapolation choice}\label{20001241}
{\small\sffamily UnsolvedMath ID 20001241 · Computer Science}
\subsection{Definitions and mathematical statement}

\paragraph{Shared notation.}Standard number systems are denoted by $\mathbb N,\mathbb Z,\mathbb Q,\mathbb R,\mathbb C$; index ranges, computational models and specialized notation follow the cited problem source and the definitions below.


For a finite design $D\subseteq\mathbb R^d$, its design ideal is $I(D)=\{f:f(v)=0\text{ for every }v\in D\}$. Polynomial functions are compared modulo $I(D)$; evaluating on $D$ identifies their restrictions to the design. A term order selects standard monomials for a Gröbner basis of $I(D)$, hence one particular basis for polynomial functions on the observed design. Different representatives agree on $D$ but can make different predictions elsewhere. The question concerns removing dependence on this arbitrary order when fitting a model.
\paragraph{Statement recorded in the supplied source.}
Question 64 (Laubenbacher). Choosing models modulo the design ideal depends on the term order chosen. How can we eliminate this choice of term order from the problem of finding a model that fits the data? Richards: Could you use random search or genetic algorithms to solve this problem? \sref{20001241}{2}
\paragraph{Scope qualification.}
The source does not fix a model-selection loss, off-design prediction rule, or uniqueness criterion.
\paragraph{Source status.}
The supplied catalogue labels this record partially solved. Its source-status metadata was checked against the supplied snapshot on 2026-10-11. This import does not independently establish current open status or mathematical correctness.
\subsection{Short English statement}
Can polynomial models fitted to a design be chosen without depending on an arbitrary monomial order?
\subsection{Sources}
\begin{itemize}
\item[\hypertarget{source:20001241:1}{S1}] Ulam AI and the UnsolvedMath Contributors, UnsolvedMath: A Curated Collection of Open Mathematics Problems, record 20001241 (AIM-COMPUTATION-0079). Catalogue curation and metadata: CC BY 4.0. The supplied source question is retained; definitions and the English summary are editorial additions.
\url{https://huggingface.co/datasets/ulamai/UnsolvedMath}.
\item[\hypertarget{source:20001241:2}{S2}] AIM Workshop Problem Lists, original problem source linked by catalogue record AIM-COMPUTATION-0079
\url{https://aimath.org/WWN/compalgstat/compalgstat.pdf}.
\end{itemize}
\paragraph{Preparation and review.}
Original question attributed to the cited source and UnsolvedMath contributors. Codex (GPT-6; exact serving revision not exposed) assisted with definitions, English summaries, source inspection and formatting. No new human mathematical authorship or independent proof verification is asserted. Canonical ID 20001241 and source code AIM-COMPUTATION-0079 checked against the complete supplied registry and existing repository statements; known flagged repetitions and obvious identical conjectures excluded. Full mathematical-equivalence/literature review has not been performed. Statement preparation and final review followed contribution prompts C and D; automated checks verify structure and identity, not open status.
\label{end:20001241}
\end{document}
